\documentclass[12pt,a4paper]{article}
\usepackage{amsmath,amssymb}
\usepackage{geometry}
\geometry{margin=2.5cm}
\title{\bf On the Connection Between the Twin Helices and the Riemann Helix}
\author{}
\date{}

\begin{document}
\maketitle

The connection between the twin conical helices (built from the Dirichlet series terms) and the Riemann helix (built from the zeros) is inevitably mediated by the Weil explicit formula.  Far from polluting the geometric significance, \textbf{the explicit formula is precisely the geometric identity that links the two}.  Its geometric restatement as the torsion of the Riemann helix being a prime‑supported Dirac comb turns a purely number‑theoretic theorem into a statement about the curvature of a space curve.  The twin helices then provide the classical dynamical system whose periodic orbits — after a M\"obius sieve — have lengths $\log p$.  In this way the explicit formula becomes the bridge between two geometric worlds, and the resulting picture is even more significant than either object alone.

\section{The Two Geometric Realms}
\begin{itemize}
\item \textbf{Additive / Dirichlet side:} The twin conical helices encode the integers $n$ as heights, with radii $n^{-\sigma}$ and phases $\pm t\log n$.  Their intersection polygon $\mathcal{P}$ extracts the real extremal points.
\item \textbf{Multiplicative / Prime side:} The Riemann helix encodes the zeros $\rho_n$ via its phase $\phi(t)$ ($\phi(\gamma_n)=2\pi n$).  Its torsion is the prime distribution $\sum_{p,m}\frac{\log p}{p^{m/2}}\delta(t-m\log p)$.
\end{itemize}
The two sides ``talk'' to each other through the Euler product $\prod_{p}(1-p^{-s})^{-1} = \sum_n n^{-s}$, which is a combinatorial identity linking integers and primes.  The analytic force of this identity is captured exactly by the Weil explicit formula.  Thus any connection between the two geometric objects must, at its core, invoke that identity.  The explicit formula is not an external additive; it \emph{is} the connection.

\section{The Explicit Formula as a Geometric Statement}
Conjecturally, the torsion of the Riemann helix is
\[
\tau_{\mathrm{HP}}(t) = \sum_{p,m}\frac{\log p}{p^{m/2}}\,\delta(t-m\log p) + \text{smooth background}.
\]
This single equation contains the entire explicit formula in a distributional form.  It can be read as:
\begin{quote}
\textit{The twist of the curve that aligns the zeros is exactly the set of point masses located at the logarithms of prime powers, weighted by the von Mangoldt function divided by $p^{m/2}$.}
\end{quote}
This is a purely differential‑geometric statement.  Its derivation from the classical explicit formula is a hard analytic problem (equivalent to the Eigencoefficient Prime Sum Conjecture), but once assumed, it expresses the link in the most vivid possible way.

\section{How the Twin Helices Fit In}
The twin helices $\mathcal{C}_\pm$ sample all integers, not just primes.  However, by taking a suitable superposition weighted by the M\"obius function $\mu(n)$, one can eliminate the composite integers and leave only the prime powers.  This is the classical identity
\[
\sum_{n} \frac{\mu(n)}{n}\log \zeta(ns) = \sum_{p} \log(1-p^{-s}),
\]
which shows that the primes are the ``primitive'' elements in the Dirichlet series.  Geometrically, the M\"obius‑weighted sum of conical helices corresponds to a curve whose classical periodic orbits are exactly the prime logarithms.  The Riemann helix is then the \emph{quantum} counterpart of that classical system; its torsion is the spectral manifestation of those prime orbits.  This is the Gutzwiller–Selberg philosophy: the trace of the quantum operator equals a sum over classical periodic orbits.  Here the operator is the CCM Dirac operator, and the periodic orbits are the M\"obius‑filtered conical helices.

Thus the twin helices naturally generate the prime orbits via M\"obius inversion, and the explicit formula turns the trace of the quantum operator into the prime sum.  The connection is therefore:
\[
\text{(twin helices)} \xrightarrow{\text{M\"obius}} \text{(prime geodesics)} \xleftrightarrow{\text{explicit formula}} \text{(torsion of Riemann helix)}.
\]

\section{Significance Is Amplified, Not Polluted}
Rather than being a pollutant, the explicit formula becomes the central geometric identity that unifies the two constructions.  The fact that we must invoke it does not diminish the geometric picture; on the contrary, it gives the explicit formula its long‑awaited geometric meaning.  The Riemann helices and the twin helices are no longer arbitrary illustrations; they are \emph{the} geometric objects that the explicit formula connects.  The resulting unity of number theory and geometry is the very significance of the approach.

\end{document}